\documentclass[11pt,a4paper]{article}

\usepackage[margin=2.5cm]{geometry}
\usepackage{amsmath,amssymb,amsthm}
\usepackage{booktabs}
\usepackage{hyperref}
\usepackage{xcolor}
\usepackage{enumitem}

\hypersetup{colorlinks=true,linkcolor=blue!50!black,urlcolor=blue!50!black}

\newcommand{\Ys}{Y}
\newcommand{\Zs}{Z^{\mathrm{s}}}
\newcommand{\Zm}{Z^{\mathrm{m}}}
\newcommand{\fs}{f_{\mathrm{s}}}
\newcommand{\fm}{f_{\mathrm{m}}}
\newcommand{\pa}{\mathrm{pa}}
\newcommand{\ch}{\mathrm{ch}}
\newcommand{\given}{\,\vert\,}
\DeclareMathOperator{\PL}{PL}

\title{The metabolite Markov random field:\\
model, estimation, and why the inference was changed}
\author{Notes accompanying the \texttt{markov-random-field} code base}
\date{Last updated: 14 August 2026}

\begin{document}
\maketitle

\begin{abstract}
This document records the probabilistic model implemented in the code base, an error in the
way its parameters were being estimated, the corrections that were made, and the reasoning
behind each. It is maintained alongside the code; section~\ref{sec:status} lists what is
implemented and what is not. A companion file \texttt{CONVERSATION.md} records the
empirical measurements that motivated each change.
\end{abstract}

\tableofcontents

%==========================================================================
\section{The model}
\label{sec:model}

\subsection{State space}

There are two trees: a species tree $T_{\mathrm{s}}$ and a molecule tree $T_{\mathrm{m}}$.
Write $S$ for the number of species \emph{leaves} and $M$ for the number of molecule leaves.

The observable latent object is the binary matrix
\[
  \Ys \in \{0,1\}^{S \times M},
  \qquad
  \Ys_{ij} = 1 \iff \text{species leaf } i \text{ produces molecule leaf } j .
\]

Each tree additionally carries latent states at its \emph{internal} nodes. For the species
tree these are
\[
  \Zs \in \{0,1\}^{I_{\mathrm{s}} \times M},
\]
where $I_{\mathrm{s}}$ is the number of internal species nodes (roots included): for every
molecule leaf $j$ there is one binary state per internal species node. Symmetrically
$\Zm \in \{0,1\}^{S \times I_{\mathrm{m}}}$.

\paragraph{Cliques.} A \emph{clique} of the species tree is the set of species-tree nodes
obtained by fixing a molecule leaf $j$; there are $M$ of them, indexed by $c = j$. Within a
clique the leaves carry $\Ys_{\cdot j}$ and the internal nodes carry $\Zs_{\cdot j}$. The
cliques of the molecule tree are indexed by species leaves. This is the meaning of the
parameter names \texttt{species\_log\_nu\_molecules\_<j>}: a species-tree parameter indexed
by a molecule leaf.

\subsection{The per-tree factor}

Along each branch of a tree the binary state evolves as a two-state continuous-time Markov
chain. For clique $c$ with switching rate $\nu_c > 0$ and stationary occupancy
$\alpha_c \in (0,1)$ the generator is
\[
  \Lambda_c \;=\; \nu_c
  \begin{pmatrix}
    -\alpha_c & \alpha_c\\[2pt]
    1-\alpha_c & -(1-\alpha_c)
  \end{pmatrix},
  \qquad
  P_c(t) \;=\; e^{\Lambda_c t}
  \;=\;
  \begin{pmatrix}
    1-\alpha_c + \alpha_c e^{-\nu_c t} & \alpha_c\bigl(1 - e^{-\nu_c t}\bigr)\\[2pt]
    (1-\alpha_c)\bigl(1 - e^{-\nu_c t}\bigr) & \alpha_c + (1-\alpha_c) e^{-\nu_c t}
  \end{pmatrix}.
\]
The stationary distribution is $(1-\alpha_c,\ \alpha_c)$, so $\alpha_c$ is the long-run
probability of state $1$ and $\nu_c$ sets how fast a branch forgets its parent.

Branch lengths are discretised onto a grid of $K$ bins. With
$\delta = 2/(K+1)$, bin $k$ represents length $\delta\,(k + \tfrac12)$, and the code builds
the bin matrices by the recursion
$P^{(0)} = e^{\Lambda_c \delta/2}$, $P^{(k)} = P^{(k-1)} e^{\Lambda_c \delta}$
(\texttt{TMatrices::\_fill\_matrices}). The bin indices are constrained so that the mean
branch length is $1$, which is what makes $\nu_c$ identifiable at all: only the product
$\nu_c t$ enters $P$.

For a single tree, writing $x_v$ for the state at node $v$, $\pa(v)$ for its parent and
$t_v$ for the length of the branch above it, the factor is
\begin{equation}
  \fs(\Ys, \Zs \given \theta_{\mathrm{s}})
  \;=\;
  \prod_{c=1}^{M}
  \Biggl[
    \pi_c\bigl(x_{r_c}\bigr)
    \prod_{v \ne r_c} P_c\bigl(t_v\bigr)_{x_{\pa(v)},\, x_v}
  \Biggr],
  \label{eq:factor}
\end{equation}
where $r_c$ is the root and $\pi_c = (1-\alpha_c, \alpha_c)$. Taken on its own,
\eqref{eq:factor} is a \emph{proper, normalised} distribution over $(\Ys, \Zs)$:
$\sum_{\Ys,\Zs} \fs = 1$.

\subsection{The joint field: a product of experts}

The two trees are combined multiplicatively, sharing the leaf layer $\Ys$:
\begin{equation}
  p\bigl(\Ys, \Zs, \Zm \given \theta\bigr)
  \;=\;
  \frac{\fs(\Ys, \Zs \given \theta_{\mathrm{s}})\,
        \fm(\Ys, \Zm \given \theta_{\mathrm{m}})}
       {C(\theta)},
  \qquad
  C(\theta) \;=\; \sum_{\Ys} p_{\mathrm{s}}(\Ys)\, p_{\mathrm{m}}(\Ys),
  \label{eq:joint}
\end{equation}
where $p_{\mathrm{s}}(\Ys) = \sum_{\Zs} \fs$ is the species tree's marginal over the leaf
layer, and likewise for $p_{\mathrm{m}}$.

\begin{quote}
\textbf{This is the crux.} Each factor is separately normalised, but their product is not:
$\Ys$ is claimed by both. The constant $C(\theta)$ is a function of every
$\alpha$, $\nu$ and branch length in \emph{both} trees, and it is a sum over $2^{SM}$
configurations --- intractable.
\end{quote}

Note that $C(\theta)$ becomes independent of $\theta_{\mathrm{s}}$ exactly when
$p_{\mathrm{m}}$ is uniform over $\Ys$; that special case is used in
section~\ref{sec:evidence} as a control.

\subsection{Observation models}

Two data sources sit above $\Ys$ (a third, mass spectrometry, is stubbed out).

\paragraph{LOTUS.} $L_{ij} \in \{0,1\}$ records whether an occurrence was reported. With
$o^{\mathrm{s}}_i = \log(1 + \text{papers}_i)$ and likewise $o^{\mathrm{m}}_j$, the
detection probability is
\[
  \rho_{ij} \;=\; \bigl(1 - e^{-\gamma_{\mathrm{s}} o^{\mathrm{s}}_i}\bigr)
                  \bigl(1 - e^{-\gamma_{\mathrm{m}} o^{\mathrm{m}}_j}\bigr),
  \qquad
  \Pr(L_{ij}=1 \given \Ys_{ij}) =
  \begin{cases}
    \rho_{ij} & \Ys_{ij} = 1,\\
    \varepsilon & \Ys_{ij} = 0 .
  \end{cases}
\]

\paragraph{Simple error model.} $D_{ij}$ reports $\Ys_{ij}$ flipped independently with
probability $\varepsilon_{\mathrm{simple}}$.

%==========================================================================
\section{What was wrong with the estimation}
\label{sec:wrong}

\subsection{The latent-field updates were correct}

$C(\theta)$ cancels in any \emph{conditional} of \eqref{eq:joint}, so the Gibbs updates for
the field were already right:
\begin{align}
  \Pr\bigl(\Ys_{ij} = y \given \text{rest}\bigr)
  &\;\propto\;
  P^{\mathrm{s}}_{j}\bigl(t_i\bigr)_{\,\Zs_{\pa(i),j},\, y}
  \cdot
  P^{\mathrm{m}}_{i}\bigl(t_j\bigr)_{\,\Zm_{i,\pa(j)},\, y}
  \cdot
  \Pr(\text{data}_{ij} \given y),
  \label{eq:ycond}\\[4pt]
  \Pr\bigl(z_v = s \given \text{rest}\bigr)
  &\;\propto\;
  P_c\bigl(t_v\bigr)_{\,z_{\pa(v)},\, s}
  \prod_{u \in \ch(v)} P_c\bigl(t_u\bigr)_{\,s,\, x_u}.
  \label{eq:zcond}
\end{align}

Two structural consequences of \eqref{eq:ycond} are worth stating explicitly.

\begin{enumerate}[label=(\alph*)]
  \item The conditional of $\Ys_{ij}$ depends on the rest of the field \emph{only} through
        the two parent states $\Zs_{\pa(i),j}$ and $\Zm_{i,\pa(j)}$. Hence
        \textbf{all $\Ys$ cells are conditionally independent given $(\Zs,\Zm)$}, and the
        existing sweep over $\Ys$ is already an \emph{exact block draw}.
  \item Given $\Ys$, the two trees decouple: $\Zs \perp \Zm \given \Ys$.
\end{enumerate}

So the sampler is a three-block Gibbs scheme on $\{\Ys\}, \{\Zs\}, \{\Zm\}$.

\subsection{The parameter update was not}

For $\theta$ the constant does \emph{not} cancel. The correct conditional is
\[
  p\bigl(\theta_{\mathrm{s}} \given \Ys, \Zs, \Zm\bigr)
  \;\propto\;
  \fs(\Ys,\Zs \given \theta_{\mathrm{s}})\;
  C(\theta)^{-1}\;
  p(\theta_{\mathrm{s}}),
\]
but \texttt{TTree::\_update\_nu\_or\_alpha} scored proposals with $\log \fs$ alone. This is
a \emph{doubly intractable} posterior being sampled as if it were an ordinary one. Two
distinct failures follow.

\paragraph{Bias.} The estimate is displaced by $\partial \log C / \partial \theta$. Measured
magnitude: $\pm 0.25$ to $0.30$ in $\log \nu$ and in $\alpha$, with the sign depending on the
operating point, vanishing entirely when the second tree is made neutral.

\paragraph{An unbounded direction.} As $\nu_c \to 0$ every $P_c(t) \to I$, so a frozen field
drives $\log \fs$ towards its maximum. In the true posterior this is paid for by $-\log C$
growing equally; with $C$ omitted the gain is free. Empirically the field term gained
$\approx 5.6\times 10^{4}$ log units between a typical field and a frozen one, and the
sampler duly ran $\log\nu \to -18$, froze $\Ys$ to all zeros, and charged the entire data
set to the error rates.

\begin{quote}
\textbf{Why priors cannot patch this.} A Normal$(0,1)$ penalty at $\log\nu = -7.5$ costs
$28$ log units; a Beta$(1,99)$ penalty at $\varepsilon = 0.14$ costs $\approx 15$. Against
$5.6\times10^{4}$ these are three orders of magnitude too small. Only a hard constraint
(pinning $\varepsilon$) stopped the collapse, which is a diagnostic, not a fix.
\end{quote}

%==========================================================================
\section{Phase 1: maximum pseudo-likelihood on internal nodes}
\label{sec:phase1}

\subsection{Formulation}

Replace the product of unnormalised factors by a product of \emph{normalised full
conditionals}. For any subset $\mathcal{V}$ of sites,
\begin{equation}
  \log \PL(\theta) \;=\; \sum_{v \in \mathcal{V}} \log p\bigl(x_v \given x_{-v},\, \theta\bigr).
  \label{eq:pl}
\end{equation}
Every term is a probability in $[0,1]$; $C(\theta)$ appears in both numerator and
denominator of each conditional and cancels exactly. Maximum pseudo-likelihood is a
composite likelihood and is consistent under the usual regularity conditions, for
\emph{any} choice of $\mathcal{V}$ --- a fact this implementation exploits.

Phase 1 takes $\mathcal{V} = \{\text{internal nodes of the tree being updated}\}$. By
\eqref{eq:zcond} those conditionals involve only that tree's own transition matrices, so no
cross-tree information is required:
\begin{equation}
  p\bigl(z_v = s \given x_{-v}\bigr)
  \;=\;
  \frac{
    P_c(t_v)_{\,z_{\pa(v)},\,s} \displaystyle\prod_{u \in \ch(v)} P_c(t_u)_{\,s,\,x_u}
  }{
    \displaystyle\sum_{s' \in \{0,1\}}
    P_c(t_v)_{\,z_{\pa(v)},\,s'} \prod_{u \in \ch(v)} P_c(t_u)_{\,s',\,x_u}
  },
  \label{eq:plterm}
\end{equation}
with $P_c(t_v)_{z_{\pa(v)},s}$ replaced by $\pi_c(s)$ at a root.

Two things are gained beyond the removal of $C$:
\begin{enumerate}[label=(\roman*)]
  \item The $\nu \to 0$ direction is no longer unbounded --- numerator and denominator move
        together.
  \item The old code scored every $\Ys$ cell \emph{twice}, once per tree, which
        over-weighted the field relative to the data by roughly a factor of two. Dropping
        the leaves removes that double counting outright.
\end{enumerate}

The cost is statistical efficiency: roughly half the nodes, so a standard error worse by
about $\sqrt2$. Since per-clique $\nu$ was already not estimable (SE $\approx 0.5$ against a
true spread $\approx 0.3$), little was lost.

\subsection{Implementation}

The denominator of \eqref{eq:plterm} is exactly the normaliser the $Z$ sampler already
computes, so the term is evaluated as a logistic of a log difference,
\[
  p\bigl(z_v = s\bigr) = \bigl[1 + \exp(\ell_{1-s} - \ell_{s})\bigr]^{-1},
  \qquad
  \ell_s = \log P_c(t_v)_{z_{\pa(v)},s} + \sum_{u\in\ch(v)} \log P_c(t_u)_{s,x_u},
\]
clamping the exponent at $700$ so that a vanishing conditional yields $\approx 10^{-305}$
rather than an exact zero (which would send the Hastings ratio to $-\infty$).

Code: \texttt{TTree::\_add\_pseudo\_likelihood\_term}, with
\texttt{TClique::calculate\_log\_prob\_node\_to\_children} exposed and given a runtime
\texttt{UseTryMatrix} flag.

\subsection{Result}

The field stopped collapsing with the error rates left free. With $\Ys$ and $Z$ pinned at
the truth the estimator became essentially unbiased ($+0.055$ and $+0.024$ in $\log\nu$,
against $\pm0.25$--$0.30$ before).

%==========================================================================
\section{Phase 2: leaf conditionals (not yet implemented)}
\label{sec:phase2}

\subsection{Why it matters}

After Phase 1, $\theta$ learns about $\nu$ \emph{only} through the latent $Z$ layer, and $Z$
is itself sampled given $\theta$. This closes a feedback loop: a slightly too small $\nu$
produces an over-smoothed $Z$, which in turn supports a small $\nu$. It is a bias in the
target distribution, so no amount of extra sampling removes it.

Measured (true $-0.446$):
\begin{center}
\begin{tabular}{lrr}
\toprule
regime & inferred \texttt{species\_mean\_log\_nu} & bias\\
\midrule
$\Ys$ and $Z$ both pinned at truth & $-0.402$ & $+0.055$\\
$\Ys$ pinned at truth, $Z$ free    & $-0.854$ & $-0.400$\\
$\Ys$ and $Z$ both free            & $-1.002$ & $-0.546$\\
\bottomrule
\end{tabular}
\end{center}
About three quarters of the residual bias is already present with a \emph{perfect} $\Ys$.

Adding the leaf terms anchors $\theta$ to $\Ys$, which the data constrains directly. The
extra terms are, from \eqref{eq:ycond},
\[
  p\bigl(\Ys_{ij} = y \given \text{rest}\bigr)
  \;=\;
  \frac{P^{\mathrm{s}}_{j}(t_i)_{\,\Zs_{\pa(i),j},\,y}\;
        P^{\mathrm{m}}_{i}(t_j)_{\,\Zm_{i,\pa(j)},\,y}}
       {\sum_{y'} P^{\mathrm{s}}_{j}(t_i)_{\,\Zs_{\pa(i),j},\,y'}\;
        P^{\mathrm{m}}_{i}(t_j)_{\,\Zm_{i,\pa(j)},\,y'}} .
\]
The other tree's factor does not depend on $\theta_{\mathrm{s}}$ in the numerator but does
not cancel from the denominator, which is precisely why it must be carried.

\subsection{Implementation status and an unresolved discrepancy}
\label{sec:phase2status}

Implemented and \textbf{on by default}; \texttt{--no\_leaf\_pseudo\_likelihood} restores the
internal-nodes-only behaviour of section~\ref{sec:phase1}.

Note first that the data likelihood does \emph{not} enter the leaf term. The pseudo-likelihood
being maximised is that of the field $p(\Ys, Z \given \theta)$; the factor
$\Pr(\text{data}\given \Ys)$ carries no $\theta$ and cancels from the Hastings ratio. Only the
other tree's parent term has to be carried, and only because it sits inside the normaliser.

Measured, with $\Ys$ and $Z$ pinned at the simulated truth:
\begin{center}
\begin{tabular}{lrr}
\toprule
$\mu_{\log\nu}$ (species) & coupled field & control: molecules tree neutral\\
                          & (true $-0.446$) & (true $-0.500$)\\
\midrule
internal nodes only (Phase 1) & $-0.402$ (bias $+0.044$) & $-0.355$ (bias $+0.146$)\\
with leaf terms (Phase 2)     & $-0.753$ (bias $-0.307$) & $-0.529$ (bias $-0.029$)\\
\bottomrule
\end{tabular}
\end{center}

On the control the leaf terms behave exactly as the theory says: the estimator improves, on
$\alpha$ as well ($0.585$ versus $0.523$ against a true $0.600$). On a genuinely coupled field
they make it worse.

A diagnostic build that forces the other-tree factor to a constant $\{\tfrac12,\tfrac12\}$ on
the coupled scenario gives bias $-0.046$, against $-0.307$ with the real factor. So the factor
is live and the leaf-term algebra is right, yet supplying the \emph{correct} factor is worse
than supplying a deliberately wrong constant --- which cannot be.

The one component the control does not exercise is the cross-tree lookup: when the other tree
is neutral its factor is $\{\tfrac12,\tfrac12\}$ regardless of which parent state is read, so a
wrong index cannot show up there. That lookup is exactly what matters once the trees couple.

\subsubsection*{The lookup is correct}

\texttt{TTree::check\_leaf\_lookup\_consistency()}, run with \texttt{--check\_leaf\_lookup},
recomputes $P_{\mathrm{other}}$ for every cell by a route that does \emph{not} assume the
clique/leaf role swap: it builds the cell's index in $\Ys$ space from (clique, leaf), and from
that derives the other tree's clique through \texttt{get\_clique(IndexArray)} --- the call the
$\Ys$ sweep itself makes --- together with its leaf node and its parent's state.

All $16384$ cells agree, for both trees. Mutation-checked, which matters because both test
scenarios are $128\times128$ and a clique/leaf swap would not even go out of range: swapping
the roles makes $16256$ of $16384$ cells disagree (the survivors are exactly the diagonal
$\text{clique}=\text{leaf}$), and swapping the two dimensions of the $Z$ index makes $5774$
disagree.

\subsubsection*{Then what causes the pinned-field discrepancy?}

Not the lookup. The likeliest remaining explanation is that the leaf terms \emph{couple the two
trees' parameter estimates}: the species leaf term contains $\alpha_{\mathrm{m}}$ and
$\nu_{\mathrm{m}}$, which are themselves weakly identified and wandering. The score for
$\theta_{\mathrm{s}}$ has zero expectation only when $\theta_{\mathrm{m}}$ sits at the truth,
so each tree inherits the other's current error --- a coupling the internal-node terms do not
have, since those involve one tree alone. This is a hypothesis, not a demonstrated fact.

In practice it is dominated by the gain: with an \emph{inferred} field the leaf terms take Y
recovery from $0.681$ to $0.920$ and bring both error rates back to their simulated values.

\subsection{The only cross-tree quantity needed}

For the species update, cell $(i,j)$ requires the single bit $\Zm_{i,\pa(j)}$. For a fixed
species clique (molecule leaf $j$) that is a fixed molecule-internal index with varying $i$
--- a strided walk through $\Zm$. Transposing those bits once per iteration into a bit-set
costs $SM$ bits ($\approx 3$\,MB at $S=M=5000$) and one $O(SM)$ pass, after which the inner
loop is a bit lookup plus one $2\times2$ matrix.

%==========================================================================
\section{Exact block sampling of \texorpdfstring{$Z$}{Z}}
\label{sec:block}

\subsection{Motivation}

$\Ys$ is already drawn as an exact block. $Z$ is not: the code visits internal nodes one at
a time using \eqref{eq:zcond}. On a strongly coupled field that produces very long
correlation times --- measured effective sample size was $\approx 10$ over $10^4$
post-burn-in iterations, i.e.\ about one independent draw per $1000$ iterations.

\subsection{The algorithm}

Conditional on $\Ys$ and $\theta$, each clique's internal states form a two-state
tree-structured hidden Markov model, which admits an \emph{exact} joint draw in time linear
in the number of nodes.

\paragraph{Upward pass (Felsenstein pruning).} For each node $v$ define the partial
likelihood of everything observed below $v$:
\[
  L_v(s) \;=\; \Pr\bigl(\Ys \text{ in the subtree below } v \given x_v = s\bigr)
  \;=\;
  \prod_{u \in \ch(v)}
  \Bigl(\textstyle\sum_{s'} P_c(t_u)_{s,s'}\, L_u(s')\Bigr),
\]
with $L_u(s') = \mathbb{1}[s' = \Ys_u]$ at a leaf, so that a leaf child contributes simply
$P_c(t_u)_{s,\,\Ys_u}$. Computed in post-order (children before parents). Working in logs,
each $L_v$ is renormalised by its maximum, which is harmless because only ratios are ever
used.

\paragraph{Downward pass (backward sampling).} At a root,
\[
  \Pr(x_{r} = s \given \Ys) \;\propto\; \pi_c(s)\, L_{r}(s),
\]
and thereafter, for an internal node $v$ whose parent has already been sampled,
\[
  \Pr\bigl(x_v = s \given x_{\pa(v)},\, \Ys\bigr) \;\propto\; P_c(t_v)_{\,x_{\pa(v)},\,s}\, L_v(s).
\]
Traversing in reverse post-order (parents before children) yields a draw from the exact
joint conditional $p(Z_{\cdot c} \given \Ys, \theta)$.

\paragraph{Cost.} $O(\text{number of edges})$ per clique --- the same order as a single
site-by-site sweep, but returning an independent draw rather than a highly correlated one.
The target distribution is unchanged: this is an exact block Gibbs step, not an
approximation.

Note the upward pass reads only \emph{leaf} states and transition matrices; it never reads
internal states. The downward pass may therefore overwrite internal states in place.

\subsection{Measured outcome: the block on \texorpdfstring{$Z$}{Z} alone does not help}
\label{sec:blockresult}

Implemented and compared against the single-site sweep on identical data and seed. Wall
clock was indistinguishable (467\,s versus 471\,s for $3\times10^4$ iterations), but:

\begin{center}
\begin{tabular}{lrr}
\toprule
 & block & single-site\\
\midrule
ESS of the field density (second half) & 9 & 9\\
iterations per independent draw        & 1710 & 1676\\
field density, second-half mean        & $-27489$ (sd 7210) & $-36660$ (sd 4344)\\
\texttt{species\_mean\_log\_nu} (true $-0.446$) & $-1.569$ & $-1.186$\\
\bottomrule
\end{tabular}
\end{center}

Two things follow. First, \textbf{the hypothesis that single-site updating of $Z$ was the
mixing bottleneck was wrong}. Second, the two schemes disagree on the posterior although both
leave $p(Z \given \Ys,\theta)$ invariant and should therefore share a stationary
distribution.

\subsection{Verification by exact enumeration}
\label{sec:enumtest}

The discrepancy above admits two readings: an implementation error, or a correct sampler
reaching regions the single-site sweep could not. These were separated by direct
verification.

The algorithm was lifted out of the storage plumbing into a templated free function
(\texttt{sample\_clique\_states\_exact}), with the tree supplied as a set of callables. The
production call site binds them to \texttt{TTree}/\texttt{TCurrentState} at no cost, while
the tests drive the same code from plain vectors.

For a clique small enough to enumerate, the target is computable in closed form:
\[
  p\bigl(Z \given \Ys\bigr)
  \;=\;
  \frac{\pi_c(x_{r})\prod_{v \ne r} P_c(t_v)_{x_{\pa(v)},\,x_v}}
       {\sum_{Z'} \pi_c(x'_{r})\prod_{v \ne r} P_c(t_v)_{x'_{\pa(v)},\,x'_v}},
\]
summing over all $2^{\lvert\text{internal}\rvert}$ configurations. \texttt{tests/TUpdate\_Z\_Tests.cpp}
compares this against $2\times10^{5}$ draws on a Monte Carlo $z$-score scale, for a balanced
tree, a high switching rate, a low switching rate, a chain, and a root-with-one-leaf case
where the root marginal must equal $\alpha$ exactly. All five pass.

The tests were then checked for power by mutation: transposing the parent-to-child lookup in
the downward pass fails four of the five, and removing the pruning recursion fails three. (The
survivor in the first case is the root-only tree, which has no non-root internal node and
therefore cannot exercise that path.)

\medskip
\noindent\textbf{Verdict.} The block sampler is correct. The disagreement in
section~\ref{sec:blockresult} is therefore the \emph{second} reading: better exploration of
the field brings the chain closer to a target that is still biased by the Phase-1
pseudo-likelihood (section~\ref{sec:phase2}). The apparently worse $\mu_{\log\nu}$ is that
bias becoming more visible, not a regression.

\subsection{The right block: \texorpdfstring{$(\Ys, Z)$}{(Y,Z)} jointly}
\label{sec:yzblock}

$\Ys$ was already an exact block and $Z$ now is too, yet nothing changed. This localises the
problem to the dependence \emph{between} the blocks: a two-block Gibbs sampler mixes slowly
precisely when its blocks are strongly dependent, and here each is close to a deterministic
function of the other.

The remedy is to enlarge the block. For a species clique $j$, condition on the \emph{other}
tree only:
\begin{equation}
  p\bigl(\Ys_{\cdot j},\, \Zs_{\cdot j} \given \Zm, \text{data}\bigr)
  \;\propto\;
  \underbrace{\pi_j(x_{r})\prod_{v\ne r} P^{\mathrm{s}}_j(t_v)_{x_{\pa(v)},x_v}}
             _{\text{species-tree factor}}
  \;\cdot\;
  \prod_{i=1}^{S}
  \underbrace{P^{\mathrm{m}}_i(t_j)_{\,\Zm_{i,\pa(j)},\,\Ys_{ij}}\;
              \Pr\bigl(\text{data}_{ij}\given \Ys_{ij}\bigr)}_{\text{leaf emission}} .
  \label{eq:yzblock}
\end{equation}

Equation~\eqref{eq:yzblock} is exactly a two-state tree-structured hidden Markov model
\emph{with emissions at the leaves}. The recursion of section~\ref{sec:block} therefore
applies unchanged, except that leaves now carry an emission likelihood instead of a fixed
state and are sampled along with the internal nodes.

Two properties make this practical:
\begin{enumerate}[label=(\roman*)]
  \item Given $\Zm$, distinct species cliques share nothing, so all $M$ of them are mutually
        independent and can be updated in parallel.
  \item Both compiled data sources factorise over cells --- LOTUS contributes
        $\Pr(L_{ij}\given \Ys_{ij})$ and the simple error model
        $\Pr(D_{ij}\given \Ys_{ij})$ --- so the emission term is a product of per-cell
        factors, as \eqref{eq:yzblock} requires.
\end{enumerate}

The resulting scheme alternates two exact block draws,
$(\Ys, \Zs)\given \Zm$ and $(\Ys, \Zm)\given \Zs$, and attacks the $\Ys$--$Z$ dependence
directly rather than the within-$Z$ dependence, which has now been measured not to matter.

\subsubsection*{The generalised recursion}

Only one change to section~\ref{sec:block} is needed. Where the leaf states were observed,
each leaf now carries an \emph{emission likelihood}
\[
  e_v(s) \;=\; \Pr\bigl(\text{everything attached to } v \given x_v = s\bigr),
\]
and the leaves are drawn along with the internal nodes. The pruning recursion becomes uniform
across node types, with $e_v \equiv 1$ wherever nothing is emitted:
\begin{equation}
  L_v(s) \;=\; e_v(s) \prod_{u \in \ch(v)}
               \Bigl(\textstyle\sum_{s'} P_c(t_u)_{s,s'}\, L_u(s')\Bigr),
  \label{eq:emitrec}
\end{equation}
so a leaf, having no children, is simply $L_v = e_v$. The downward pass is unchanged except
that it now visits every node, drawing a leaf from
$P_c(t_v)_{\,x_{\pa(v)},\,s}\, e_v(s)$.

Comparing \eqref{eq:emitrec} with \eqref{eq:yzblock}, the emission for species leaf $i$ in
clique $j$ is
\[
  e_i(s) \;=\;
  \underbrace{P^{\mathrm{m}}_i(t_j)_{\,\Zm_{i,\pa(j)},\,s}}_{\text{other tree}}
  \;\cdot\;
  \underbrace{\Pr\bigl(L_{ij}\given s\bigr)\,\Pr\bigl(D_{ij}\given s\bigr)}_{\text{data}} .
\]
That both data sources factorise over cells is precisely what makes this an emission, rather
than a term coupling the leaves to one another.

\paragraph{Status.} Implemented as \texttt{sample\_clique\_states\_with\_emissions} and
verified the same way as section~\ref{sec:enumtest}, but now enumerating the \emph{joint}
$2^{\lvert\text{nodes}\rvert}$ space rather than only the internal nodes: a balanced tree with
pseudo-random emissions, a low switching rate, a chain, and a bridge test showing that
emissions sharp enough to pin the leaves reproduce the Z-only posterior after marginalisation.
All pass; dropping the emission from the pruning pass fails all four. It is \emph{not} yet the
production sampler --- see section~\ref{sec:wiring}.

\subsection{What wiring it in still requires}
\label{sec:wiring}

\begin{enumerate}[label=(\arabic*)]
  \item A post-order over \emph{all} nodes. The cached order currently filters the leaves out.
  \item The other tree's parent-state bit per cell, $\Zm_{i,\pa(j)}$ --- a strided walk
        through $\Zm$, cured by the transposed bit-set of section~\ref{sec:phase2}.
  \item Per-cell data likelihoods along a \emph{column} of $\Ys$. The two numbers needed are
        already produced by \texttt{TLotus::calculate\_LL\_update\_Y} and
        \texttt{TSimpleErrorModel::probabilities\_for\_Y\_update}, but both are fed from a
        buffer filled along the last dimension, i.e.\ across the grain of the species pass.
  \item Restructuring \texttt{TMarkovField}: today it sweeps $\Ys$ cell by cell and then $Z$
        per tree; the joint block replaces both.
\end{enumerate}

Item (3) is shared with Phase 2 (section~\ref{sec:phase2}), so the two are best built
together.

%==========================================================================
\section{Evidence and controls}
\label{sec:evidence}

The claims above rest on a small number of controlled experiments; they are recorded in
full in \texttt{CONVERSATION.md}. The most load-bearing:

\begin{description}[style=nextline]
  \item[Neutral-tree control] Setting the second tree to $\alpha = 0.5$ and $\nu$ above the
    stationary threshold makes $p_{\mathrm{m}}$ uniform, hence $C$ independent of
    $\theta_{\mathrm{s}}$. The bias in $\log\nu$ went from $+0.29$ to $-0.012$ and in
    $\alpha$ from $+0.25$ to $-0.004$. This is what identifies $C(\theta)$ as the culprit
    rather than, say, finite-sample bias.
  \item[Burn-in control] Re-simulating the field with $15\times$ more iterations changed the
    bias by $0.001$, excluding simulation non-convergence.
  \item[MLE control] A parametric bootstrap on a single-tree CTMC gave a finite-sample bias
    of $-0.044$, an order of magnitude too small to explain the observation.
  \item[Pinning ladder] Pinning $\Ys$ and $Z$, then only $\Ys$, then neither, localises the
    residual bias to the $Z$ feedback loop (section~\ref{sec:phase2}).
\end{description}

%==========================================================================
\section{Identifiability notes}
\label{sec:ident}

\paragraph{Per-clique $\nu$.} From $n$ branches the standard error of $\widehat{\log\nu}$ is
$\approx 0.53$ at $n = 254$. Simulations with a true spread of $\mathrm{sd} \approx 0.09$
therefore cannot be recovered: the attainable correlation between truth and estimate is
$\approx 0.17$. This is a property of the design, not a defect of the code, and it argues
for sharing $\nu$ across cliques (or making it a function of covariates) rather than
estimating $M$ separate values.

\paragraph{$\gamma$.} Detection saturates: with $\log(1+\text{papers})$ of order $1.6$, the
factor $1 - e^{-\gamma o}$ exceeds $0.997$ once $\gamma \gtrsim 5$, so the likelihood is
flat above that. Simulating at $\gamma = 5$ places the truth inside the flat region. A
Gamma$(2,4.6)$ prior has mean $\approx 0.43$, which puts detection at $\tfrac12$ for the
median paper count and negligible mass on the saturated corner.

\paragraph{$\nu$ versus branch lengths.} Only the product $\nu_c t$ enters $P$. Identifiability
comes solely from the constraint that the mean binned branch length is $1$.

%==========================================================================
\section{Simulation-side corrections}
\label{sec:sim}

\paragraph{Variance versus standard deviation.} \texttt{model\_validation/src/params.py} drew
$\log\nu_c \sim \mathcal{N}(\mu, \texttt{var\_log\_nu})$ by passing the variance to
\texttt{numpy}'s \texttt{scale} argument, which is the \emph{standard deviation}. The C++
prior uses $\mathrm{sd} = \sqrt{\mathrm{var}}$. The simulated spread was therefore the square
root of the intended one, and the reported ``true'' \texttt{var\_log\_nu} was actually a
standard deviation. Corrected to \texttt{np.sqrt(...)}.

\paragraph{Parameters are drawn from their priors when simulating.} Any parameter not fixed
from a file is drawn from its prior during \texttt{simulate}, not set to the CLI default.
Changing a prior therefore changes what gets simulated --- which is desirable (one then
simulates from the prior one infers under) but shifts the recorded truths between runs.

%==========================================================================
\section{Status}
\label{sec:status}

\begin{center}
\begin{tabular}{llp{6.2cm}}
\toprule
Change & Status & Note\\
\midrule
\texttt{sqrt} in \texttt{params.py}      & done & Simulator only.\\
Informative priors on $\gamma$, $\varepsilon$, $\mu_{\log\nu}$ & done &
  Fixes $\gamma$; the other two are defensible but do not stop the collapse.\\
Phase 1 pseudo-likelihood (internal nodes) & done & Stops the collapse; removes the
  coupling bias given a known field.\\
Block sampling of $Z$                    & done, \textbf{verified} &
  Correct (section~\ref{sec:enumtest}), but does not improve mixing. Kept: it is exact and
  free. The old sweep remains available behind \texttt{--single\_site\_Z}.\\
Joint $(\Ys,Z)$ clique block             & recursion done \& verified, \textbf{not wired} &
  Section~\ref{sec:yzblock}. The measurement in section~\ref{sec:blockresult} says this, not
  the $Z$-only block, is where the mixing win is. Remaining work in
  section~\ref{sec:wiring}.\\
Phase 2 leaf conditionals                & done, \textbf{on by default} &
  Cross-tree lookup verified cell by cell. Takes $\Ys$ recovery from $0.68$ to $0.92$ on the
  full pipeline. One unexplained oddity in the pinned-field regime;
  section~\ref{sec:phase2status}.\\
Shared (non per-clique) $\nu, \alpha$    & proposed & Modelling decision;
  section~\ref{sec:ident}.\\
\bottomrule
\end{tabular}
\end{center}

\end{document}
